\section{Why separate medial transitivity does not give one test}
The common-quotient condition in the preceding section cannot simply be
replaced by medial-section transitivity of the two inputs separately.
We give two isomorphic polytopes with different compatible rootings that
satisfy all the local and symmetry hypotheses in [CH, Corollaries~6.4,
6.6, Theorem~6.7 and Proposition~6.8]. The indicated base test passes,
although their mix is not polytopal. In particular, this is not a
counterexample depending on nonuniform facets or on a nonpolytopal input.

\subsection{A six-coordinate action}
Let $S=S_5$, with $s_i=(i+1\ i+2)$ for $0\leq i\leq3$.
Its six Sylow $5$-subgroups are indexed $0,\ldots,5$ using the respective
cycle generators
\[
 (12345),\ (12354),\ (12453),\ (12435),\ (12543),\ (12534).
\]
Conjugation induces a permutation action $\pi$ on these six coordinates.
Writing each image list in order $0,\ldots,5$, the four permutations are
\[
\begin{array}{c|rrrrrr}
i&\pi(s_i)(0)&\pi(s_i)(1)&\pi(s_i)(2)&\pi(s_i)(3)&\pi(s_i)(4)&\pi(s_i)(5)\\\hline
0&4&2&1&5&0&3\\
1&5&3&4&1&2&0\\
2&3&2&1&0&5&4\\
3&1&0&4&5&2&3
\end{array}                                                    \tag{S1}
\]
Use the induced left permutation action on
$V=\{v\in\mathbb F_3^6:\sum_hv_h=0\}$, so $(gv)_h=v_{\pi(g)^{-1}h}$.
In $V\rtimes S$ put $R_i=(u_i,s_i)$, where
\[
 u_0=u_2=u_3=0,\qquad u_1=(-1,-1,0,1,0,1).
\]
Multiplication is $(u,g)(v,h)=(u+gv,gh)$. Direct substitution in (S1)
shows $R_i^2=1$ and $[R_i,R_j]=1$ for $|i-j|>1$; also
\[
 (R_0R_1)^3=(R_2R_3)^3=1,\qquad
 c:=(R_1R_2)^3=(v,1),\quad v=(1,1,-1,-1,1,-1).       \tag{S2}
\]
These identities are small coordinate calculations, all displayed inputs
being fixed elements of $\mathbb F_3$ and $S_5$.

The four marks generate all of $G=V\rtimes S$. Indeed their projection
is $S$, and their normal closure of $c$ contains the $S$-span of $v$.
The five vectors $v,s_0v,s_3v,s_0s_3v,s_1s_0v$ have rank five over
$\mathbb F_3$; their rows are
\[
\begin{pmatrix}
1&1&2&2&1&2\\1&2&1&2&1&2\\1&1&1&2&2&2\\
2&1&1&2&1&2\\2&2&1&2&1&1
\end{pmatrix}.
\]
They lie in $V$, so span $V$. Consequently $|G|=3^5\cdot120=29160$.

For $h\in\{0,\ldots,5\}$ let $N_h=\{v\in V:v_h=0\}$ and form the
colored Schreier graph $P_h=N_h\backslash G$, with adjacency given by
right multiplication by $R_i$. Using right rather than left connection
actions merely reverses the names of words; the colored path statements
are unchanged. A flag can be written $(g,k)\in S\times\mathbb F_3$, and
its $i$-adjacency is
\[
 (g,k)^i=(gs_i,k+(g u_i)_h).                            \tag{S3}
\]
Thus every $P_h$ has $360$ flags and covers the simplex flag graph
$T=S$. The cover is regular with deck group $V/N_h\cong C_3$.
Equivalently, the preimages of $N_h\trianglelefteq V$ in the universal
connection group satisfy exactly the source's definition of
$T$-admissibility. Since $T$ is a maniplex, so are these covers.

\subsection{Every input is a medial-section-transitive polytope}
Write $A=\gen{R_0,R_1,R_2}$, $B=\gen{R_1,R_2,R_3}$ and
$C=\gen{R_1,R_2}$. Their translation kernels under $G\to S$ are
\[
\begin{aligned}
 V_A&=\gen{e_0-e_5,e_1-e_2,e_3-e_4}_{\mathbb F_3},\\
 V_B&=\gen{e_0-e_3,e_1-e_5,e_2-e_4}_{\mathbb F_3},\\
 V_C&=\gen{v}_{\mathbb F_3}.
\end{aligned} \tag{S4}
\]
Here is a full justification of these kernel formulas. The projected
end groups are $S_4$, with their usual adjacent-transposition
presentations. All its Coxeter relations already hold in each lifted
end except the middle cube in (S2). Therefore its kernel is generated
by the conjugates of that cube, and equals the linear span of the
corresponding orbit of $v$. The first space in (S4) is invariant under
$s_0,s_1,s_2$, contains $v$, and the three vectors $v,s_0v,s_1s_0v$
are independent in it. The second space is invariant under
$s_1,s_2,s_3$, and contains the independent vectors $v,s_3v,s_2s_3v$.
This proves both equalities. The last follows also from the presentation
of a group generated by two involutions: $R_1R_2$ has order nine and
its projection has order three. We retain here only the elementary
Coxeter presentations of symmetric groups; the translation kernels are
computed in full above.

The restricted groups on colors $\{0,1\}$ and $\{2,3\}$ project
isomorphically to $S_3$, by (S2). Consequently each facet of $P_h$ covers
a tetrahedron with trivial fiber on its $\{0,1\}$ sections. The common
fiber criterion gives the rank-three intersection property: one of the
two fibers being intersected is trivial. Its two-color sections are
polygons, so every facet is polytopal. The dual argument proves the same
for every vertex-figure. Each facet and vertex-figure has $24\cdot3=72$
flags, since every coordinate projection of both spaces in (S4) is
surjective. Every medial section has $6\cdot3=18$ flags, since all
coordinates of $v$ are nonzero. At every flag of $P_h$ the medial fiber
over $T$ is therefore the whole $C_3$ fiber. The common criterion now
proves that $P_h$ itself is polytopal. This proves all-flag assertions,
not only a test at a selected root.

The stabilizer $H_h$ of the $h$th Sylow subgroup in $S_5$ is its
normalizer, an affine group of order $20$ on five points. To see this
directly, identify the five points with $\mathbb F_5$: translations
normalize the chosen regular $C_5$, as do the four nonzero scalings,
and a normalizer is determined by the image of one point and of a
cycle generator. This group is sharply two-transitive. Since
$V\rtimes H_h$ normalizes $N_h$, its left action on $N_h\backslash G$
commutes with all colored adjacencies. On $T$, a middle-color component
is specified by the ordered pair $(g(1),g(5))$. There are twenty such
pairs, and $H_h$ is transitive on them. Each of these components has
one connected preimage in $P_h$, as just proved. Hence $P_h$ is
medial-section-transitive. In particular its facets, vertex-figures and
medial sections have uniform isomorphism types. The medial components themselves form one
automorphism orbit, exactly as required by the source. The six inputs
are isomorphic: conjugation by a suitable $(0,a)$ sends $N_h$ to $N_k$,
and left multiplication by $(0,a)$ then gives a color-preserving map
$N_h\backslash G\to N_k\backslash G$.

\subsection{One base passes and another fails}
Consider $P_0$ and $P_1$, rooted at compatible flags specified below.
Their mix is $(N_0\cap N_1)\backslash G$, with $1080$ flags. Its flags
are $(g,k,l)$, and (S3) holds in both sheet coordinates. All its facets
and vertex-figures are polytopal, since they are mixes of polyhedra.

At a flag above $g\in S$, the three fibers of this mix over $T$ are the
images of $gV_A,gV_B,gV_C$ under the two coordinate projections
$(v_0,v_1)$. This follows by traversing a word whose projection returns
to $g$ in the corresponding color subgroup. At $g=1$, (S4) gives
\[
 D_A=D_B=\mathbb F_3^2,\qquad D_C=\gen{(1,1)}.
\]
Thus the end intersection fails: it has $54$ flags whereas the middle
component has $18$. There is also an explicit path witness. Put
$h=R_2R_3$ and use the translation $c$ from (S2). The words
\[
 (cR_0cR_0)^2\in A,\qquad (chc^{-1}h^{-1})^2\in B
\]
both carry $(1,0,0)$ to $(1,1,0)$, by (S1)--(S2). No middle-color word
can do so, since the middle translation fiber is the diagonal line.
The mix is consequently not polytopal.

Now set $g_*=s_1s_2s_3=(2\,3\,4\,5)$, and use base flags with
$(g,k,l)=(g_*,0,0)$. From (S1), $\pi(g_*)=(4,1,0,2,3,5)$.
The coordinate pair $(0,1)$ of $g_*z$ is therefore $(z_2,z_1)$.
Equations (S4) give at this base
\[
 D_A=\gen{(-1,1)},\qquad D_B=\mathbb F_3^2,
 \qquad D_C=\gen{(-1,1)}.                              \tag{S5}
\]
The base intersection thus holds. To express the test exactly as a mix
of $P_0$ with $P_1$, rather than a cover of $T$, require that the first
sheet remain fixed. Intersect the three fibers in (S5) with
$\{0\}\times\mathbb F_3$. The facet and middle fibers become trivial;
the vertex fiber is $C_3$. In particular the two end variance groups
have trivial intersection at the specified base.

Here and throughout this paper, medial-section transitivity means
transitivity on the middle-color components, as used in the proof of
[CH, Proposition~2.5] and its two-orbit tests. The parenthetical chain
type on the source's physical page~6 should instead be $\{0,n-1\}$;
it prints the complementary type. We do not silently conflate these
notions.

Both inputs are polytopes and individually medial-section-transitive,
both are admissible over the same regular $T$, all their local types
are uniform, and even every mixed end is polytopal. Nevertheless the
base conditions hold while the mix fails elsewhere. This disproves the
single-base conclusions of [CH, Corollary~6.4 and Theorem~6.7] as
printed, and of Corollary~6.6 and Proposition~6.8 under the operational
medial-section convention just specified. The source's all-flags Theorem~6.3
is unaffected. A correct replacement either keeps every flag, assumes
medial-section transitivity of the mix itself, or uses the explicitly
proved common-quotient condition $W=CV$. The same argument identifies
exactly what went wrong: separate automorphisms of the two factors need
not act together on their chosen mix.
